\chapter{Vlastní práce}
\label{chap:practical}

% ======================================================================
% Macros of the formal model
% ======================================================================
\newcommand{\lxss}{\lambda_{\mathsf{XSS}}}
\newcommand{\Lat}{\mathcal{L}}
\newcommand{\Ctxs}{\mathcal{C}}
\newcommand{\tStr}{\mathsf{Str}}
\newcommand{\tBool}{\mathsf{Bool}}
\newcommand{\tUnit}{\mathsf{Unit}}
\newcommand{\mSrc}{\mathsf{src}}
\newcommand{\mVal}{\mathsf{val}}
\newcommand{\mOut}{\mathsf{out}}
\newcommand{\mTpl}{\mathsf{tpl}}
\newcommand{\mErr}{\mathsf{err}}
\newcommand{\Inert}{\mathit{Inert}}
\newcommand{\Exec}{\mathit{Exec}}
\newcommand{\Marked}{\mathit{Mark}}
\newcommand{\CtxOf}{\mathit{ctx}}
\newcommand{\Prs}{\mathit{parse}}
\newcommand{\San}{\mathit{San}}
\newcommand{\Skel}{\mathit{skel}}
\newcommand{\Prov}{\mathcal{I}}

Předchozí kapitola shromáždila aparát, o~který se konstruktivní prevence XSS opírá: Denningovu bezpečnostní mřížku, bezpečnostní typové systémy odvozené od~$\lambda^{SEC}$ a~návrhový vzor bezpečných hodnot doložený standardem Trusted Types. Všechny tyto stavební kameny jsou převzaté. Tato kapitola z~nich staví vlastní model,~který je formuluje jako jediný celek a~doplňuje o~to, co v~převzatém aparátu chybí: o~sémantiku,~která bezpečnostním štítkům dává obsah odvozený od chování prohlížeče, a~o~důkaz, že typová pravidla tuto sémantiku respektují.

\section{Formální model toku nedůvěryhodných dat}
\label{sec:model}

\subsection{Konstrukční požadavky a~přehled modelu}
\label{sec:model-pozadavky}

Model není samoúčelný. Má sloužit jako podklad pro knihovnu,~kterou zavádí dílčí cíl~2, a~pro auditní pravidlo z~dílčího cíle~4. Z~toho plyne pět konstrukčních požadavků. Zároveň vysvětlují, proč má model právě tu podobu,~ve které je níže předložen.

\begin{description}
    \item[R1: oddělení validace od sanitace.] Dílčí cíl~1 žádá, aby model \emph{formálně zachycoval krok validace schématu i~krok sanitace}. Nestačí tedy, aby oba kroky posouvaly hodnotu stejnou mřížkou. Každý z~nich musí měnit jinou vlastnost a~model to musí umět vyjádřit.
    \item[R2: přípustnost legitimního výstupu.] Smyslem sledované aplikace je vypsat uživatelský komentář. Model,~který by takový program zamítl, je nepoužitelný. Bezpečnostní vlastnost proto nesmí znít „vstup útočníka neovlivní výstup“, ale musí přesně vymezit, \emph{v~čem} ho ovlivnit nesmí.
    \item[R3: záměna kontextu.] Nejčastější reálnou chybou není zapomenutá sanitace, ale sanitace pro nesprávný kontext: hodnota escapovaná pro tělo HTML se vloží do hodnoty atributu \parencite{samuel2011context}. Model musí tento případ odmítnout, a~to bez zvláštního pravidla ušitého na míru.
    \item[R4: přenositelnost do TypeScriptu.] Záruka má být vynucena překladačem, ne ideálním kalkulem. Struktura štítků proto musí mít protějšek v~operacích,~které TypeScript skutečně má.
    \item[R5: vymezení auditní plochy.] Model musí umět odpovědět na otázku,~kterou sekce~\ref{sec:audit} označila za jádro auditovatelnosti: \emph{kdo je oprávněn vyrobit bezpečnou hodnotu?} Odpověď je zároveň měřenou veličinou v~hodnocení dle rámce SPACE.
\end{description}

Model se skládá ze šesti vrstev,~jejichž vzájemnou závislost zachycuje schéma~\ref{fig:model-prehled}. Postup je záměrně veden zdola: nejprve se určí, \emph{co štítek znamená} v~pojmech chování prohlížeče, a~teprve pak se odvodí pravidla,~která s~ním pracují. Opačné pořadí, tedy nejprve pravidla a~pak jejich výklad, vede k~modelu, u~kterého nelze rozhodnout, zda je zvolený operátor spojení správný, či pouze pravděpodobný.

\begin{figure}[htbp]
    \centering
    \begin{tikzpicture}[
        every node/.style={font=\footnotesize},
        blk/.style={draw, rounded corners=2pt, minimum width=33mm, minimum height=14mm, align=center, thick},
        fa/.style={blk, fill=blue!8,    draw=blue!55!black},
        fb/.style={blk, fill=orange!12, draw=orange!65!black},
        fc/.style={blk, fill=green!10,  draw=green!45!black},
        fd/.style={blk, fill=gray!12,   draw=gray!60!black},
        e/.style={->, very thick, >=Stealth, gray!55!black},
        num/.style={font=\scriptsize\bfseries, gray!60!black}
    ]
        \node[fa] (lat)  at (0,    2.9) {mřížka $\Lat$\\[1pt]\scriptsize provenience $\times$ inertnost};
        \node[fa] (sem)  at (4.55, 2.9) {význam štítků\\[1pt]\scriptsize $\Exec_c$, $\Inert_c$};
        \node[fb] (lang) at (9.1,  2.9) {jazyk $\lxss$\\[1pt]\scriptsize syntaxe a~sémantika};
        \node[fb] (typ)  at (9.1,  0.4) {typová pravidla\\[1pt]\scriptsize $\Gamma \vdash e : \tau$};
        \node[fc] (thm)  at (4.55, 0.4) {korektnost\\[1pt]\scriptsize L1--L4, V1--V5};
        \node[fd] (ts)   at (0,    0.4) {realizace v~TS\\[1pt]\scriptsize vnoření $\llbracket \cdot \rrbracket$};

        \draw[e] (lat)  -- node[above, num] {4.1.2} (sem);
        \draw[e] (sem)  -- node[above, num] {4.1.3} (lang);
        \draw[e] (lang) -- node[right, num] {4.1.5} (typ);
        \draw[e] (typ)  -- node[above, num] {4.1.6} (thm);
        \draw[e] (thm)  -- node[above, num] {4.1.8} (ts);

        \draw[e, dashed, gray!45!black] (sem) -- node[right=1mm, num, align=center] {sémantická\\opora vět} (thm);
    \end{tikzpicture}
    \caption{Vrstvy modelu a~jejich závislosti. Čísla u~hran odkazují na podsekci,~která danou vrstvu zavádí. Plné šipky vyjadřují konstrukční postup, přerušovaná šipka připomíná, že věty o~bezpečnosti nejsou tvrzeními o~symbolech, ale o~sémantice štítků zavedené v~podsekci~\ref{sec:model-semantika}.}
    \label{fig:model-prehled}
\end{figure}

\subsection{Bezpečnostní mřížka: provenience a~kontextová inertnost}
\label{sec:model-mrizka}

Teoretická kapitola pracovala se dvěma mřížkami odděleně: s~lineární mřížkou integrity $\bot \sqsubseteq \mathit{mid} \sqsubseteq \top$ (schéma~\ref{fig:pipeline}) a~s~mřížkou kontextově specifických typů (diagram~\ref{fig:mrizka-kontextova}). Požadavek~R1 velí obě spojit, nikoliv však sloučit. Jde totiž o~dvě nezávislé vlastnosti stejné hodnoty.

\begin{definice}[Osa provenience]
\label{def:provenience}
Označme $\Prov = \{\bot, \mathit{mid}, \top\}$ s~uspořádáním $\bot \sqsubseteq \mathit{mid} \sqsubseteq \top$, kde
$\bot$ značí hodnotu pocházející výhradně z~textu programu,
$\mathit{mid}$ hodnotu, jejíž \emph{struktura} byla ověřena schématem, a
$\top$ hodnotu plně řízenou útočníkem.
\end{definice}

\begin{definice}[Osa kontextů]
\label{def:kontexty}
Zápis $\Ctxs = \{\mathit{html},\, \mathit{script},\, \mathit{url},\, \mathit{res}\}$ značí konečnou množinu \emph{výstupních kontextů}. Odpovídají jim bezpečné typy \texttt{SafeHtml}, \texttt{SafeScript}, \texttt{SafeUrl} a~\texttt{TrustedResourceUrl} z~dílčího cíle~2. Podmnožina $K \subseteq \Ctxs$ vyjadřuje, ve~kterých kontextech je hodnota \emph{inertní} ve smyslu definice~\ref{def:inertnost}.
\end{definice}

Zásadní je, že obě osy nesou uspořádání v~opačném smyslu. Vyšší provenience znamená \emph{méně} důvěry, zatímco větší množina kontextů znamená \emph{více} bezpečí. Bezpečnostní štítek je proto prvkem součinu,~ve kterém je druhá složka uspořádána obráceně.

\begin{definice}[Bezpečnostní mřížka]
\label{def:mrizka}
Bezpečnostní štítek je dvojice $\ell = (i, K) \in \Prov \times \mathcal{P}(\Ctxs)$. Na množině štítků definujeme
\begin{align}
    (i, K) \sqsubseteq (i', K') &\;\;\overset{\text{def}}{\Longleftrightarrow}\;\; i \sqsubseteq_{\Prov} i' \;\wedge\; K \supseteq K',
    \label{eq:m-order}\\
    (i, K) \sqcup (i', K') &\;\;=\;\; \bigl(\, i \sqcup_{\Prov} i',\;\; K \cap K' \,\bigr),
    \label{eq:m-join}\\
    (i, K) \sqcap (i', K') &\;\;=\;\; \bigl(\, i \sqcap_{\Prov} i',\;\; K \cup K' \,\bigr).
    \label{eq:m-meet}
\end{align}
Vzniklá struktura se značí $\Lat = \bigl(\Prov \times \mathcal{P}(\Ctxs),\, \sqsubseteq,\, \sqcup\bigr)$. Její nejmenší prvek je $\bot_{\Lat} = (\bot, \Ctxs)$ a~největší $\top_{\Lat} = (\top, \emptyset)$.
\end{definice}

\begin{lemma}[Korektnost mřížky]
\label{lem:mrizka}
Struktura $\Lat$ je konečný svaz a~splňuje Denningovy axiomy \ref{ax:finitness} až \ref{ax:join} ze sekce~\ref{eq:lattice}.
\end{lemma}

\begin{proof}
$\Prov$ je konečný řetězec o~třech prvcích, tedy svaz. $\bigl(\mathcal{P}(\Ctxs), \supseteq\bigr)$ je duál Booleovy algebry podmnožin konečné množiny, tedy také konečný svaz,~ve kterém supremum odpovídá průniku a~infimum sjednocení. Kartézský součin dvou svazů s~po složkách definovaným uspořádáním je svaz a~supremum se počítá po složkách. Rovnice~\eqref{eq:m-order} až \eqref{eq:m-meet} tuto konstrukci pouze rozepisují. Odtud:
\ref{ax:finitness} plyne z~$|\Lat| = 3 \cdot 2^{|\Ctxs|} = 48$;
\ref{ax:order} z~toho, že součin částečných uspořádání je částečné uspořádání;
\ref{ax:bottom} z~existence nejmenších prvků obou složek, tedy $\bot_{\Lat} = (\bot, \Ctxs)$;
\ref{ax:join} z~existence suprem v~obou složkách a~z~jejich jednoznačnosti ve svazu.
\end{proof}

Součinovou strukturu znázorňuje diagram~\ref{fig:mrizka-soucin}. Model tímto krokem \emph{zpřesňuje} tvrzení uvedené v~popisu diagramu~\ref{fig:join-action}: spojení dvou hodnot bezpečných v~různých kontextech nedává $\top_{\Lat}$, ale $(\bot, \emptyset)$. Výsledek je sice neinertní ve všech kontextech, jeho provenience však zůstává důvěryhodná. Toto rozlišení není akademické. Právě ono umožňuje v~pravidle~\textsc{T-Out} připustit programátorem zapsaný literál obsahující značky HTML.

\begin{figure}[htbp]
    \centering
    \begin{tikzpicture}[
        every node/.style={font=\scriptsize},
        n/.style={draw, circle, inner sep=1.2pt, minimum size=5.5mm, fill=white, font=\tiny},
        b/.style={draw, rounded corners=2pt, inner sep=2.5pt, align=center, font=\tiny, fill=white},
        e/.style={-, thick, gray!65!black},
        hi/.style={draw=red!70!black, thick, fill=red!8},
        lo/.style={draw=green!50!black, thick, fill=green!10},
        pan/.style={font=\footnotesize\itshape}
    ]
        % ---------- (a) provenance axis ----------
        \node[pan] at (0, 4.5) {(a) provenience $\Prov$};
        \node[n] (i0) at (0, 0.6) {$\bot$};
        \node[n] (i1) at (0, 2.0) {$\mathit{mid}$};
        \node[n] (i2) at (0, 3.4) {$\top$};
        \draw[e] (i0) -- (i1) -- (i2);
        \node[right=1.5mm of i0, font=\tiny] {text programu};
        \node[right=1.5mm of i1, font=\tiny] {ověřená struktura};
        \node[right=1.5mm of i2, font=\tiny] {vstup útočníka};

        % ---------- (b) context axis ----------
        \begin{scope}[xshift=44mm]
            \node[pan] at (0, 4.5) {(b) kontexty $\mathcal{P}(\Ctxs)^{op}$};
            \node[n] (k0) at ( 0,   0.6) {$\Ctxs$};
            \node[n] (kh) at (-1.2, 2.0) {$\{h\}$};
            \node[n] (ks) at ( 1.2, 2.0) {$\{s\}$};
            \node[n] (ke) at ( 0,   3.4) {$\emptyset$};
            \draw[e] (k0) -- (kh); \draw[e] (k0) -- (ks);
            \draw[e] (kh) -- (ke); \draw[e] (ks) -- (ke);
            \node[below=1.5mm of k0, font=\tiny] {inertní všude};
            \node[above=1.5mm of ke, font=\tiny] {inertní nikde};
        \end{scope}

        % ---------- (c) product ----------
        \begin{scope}[xshift=88mm]
            \node[pan] at (1.95, 4.5) {(c) součin $\Lat = \Prov \times \mathcal{P}(\Ctxs)^{op}$};
            \foreach \r/\y in {0/0.6, 1/2.0, 2/3.4}
              \foreach \c/\x in {0/0, 1/1.95, 2/3.9}
                \coordinate (g\r\c) at (\x, \y);

            \node[b, lo] (p00) at (g00) {$(\bot,\Ctxs)$};
            \node[b]     (p01) at (g01) {$(\bot,\{h\})$};
            \node[b]     (p02) at (g02) {$(\bot,\emptyset)$};
            \node[b]     (p10) at (g10) {$(\mathit{mid},\Ctxs)$};
            \node[b]     (p11) at (g11) {$(\mathit{mid},\{h\})$};
            \node[b]     (p12) at (g12) {$(\mathit{mid},\emptyset)$};
            \node[b]     (p20) at (g20) {$(\top,\Ctxs)$};
            \node[b]     (p21) at (g21) {$(\top,\{h\})$};
            \node[b, hi] (p22) at (g22) {$(\top,\emptyset)$};

            \foreach \r in {0,1,2} {
              \draw[e] (p\r0) -- (p\r1); \draw[e] (p\r1) -- (p\r2);
            }
            \foreach \c in {0,1,2} {
              \draw[e] (p0\c) -- (p1\c); \draw[e] (p1\c) -- (p2\c);
            }
            \node[below=1mm of p00, font=\tiny, green!45!black] {$\bot_{\Lat}$};
            \node[above=1mm of p22, font=\tiny, red!70!black] {$\top_{\Lat}$};
        \end{scope}
    \end{tikzpicture}
    \caption{Bezpečnostní mřížka jako součin dvou os. Panel~(a) zachycuje osu provenience, panel~(b) osu kontextů pro dvouprvkovou množinu $\{h, s\}$ (v~modelu jde o~čtyřprvkovou množinu, tedy o~čtyřrozměrnou krychli o~šestnácti vrcholech). Panel~(c) ukazuje výřez součinu: směrem vzhůru roste kontaminace, směrem vpravo ubývá inertních kontextů. Obě cesty vedou vzhůru v~uspořádání~$\sqsubseteq$. Hodnota \texttt{SafeHtml} odpovídá sloupci $\{h\}$ nezávisle na řádku. Právě tato nezávislost obou os je jádrem modelu.}
    \label{fig:mrizka-soucin}
\end{figure}
